
\subsection*{1.3}
\phantomsection
\label{sga3:VIIB:1.3}
\addcontentsline{toc}{subsection}{1.3. Monomorphisms and topological flatness}

\noindent\textbf{Proposition 1.3.}
\phantomsection
\label{sga3:VIIB:1.3.prop}
{\itshape Let \(f:X\to Y\) be a morphism of formal varieties over
\(k\), let \(A\) and \(B\) be the affine algebras of \(X\) and \(Y\),
and let \(g:B\to A\) be the morphism induced by \(f\). Then \(f\) is a
monomorphism in \(\operatorname{Vaf}/k\) if and only if \(g\) is
surjective.}

By \SGARefLink{sga3:VIIB:1.1}{1.1},\footnote{The editors expanded the beginning of the proof.}
the functor \(A\mapsto\operatorname{Spf}(A)\) is an anti-equivalence
from \(\operatorname{Alp}/k\) to \(\operatorname{Vaf}/k\).
Consequently, \(f\) is a monomorphism if and only if \(g\) is an
epimorphism, which it certainly is when \(g\) is surjective.

Conversely, suppose that \(g\) is an epimorphism. For every
\(\mathfrak n\in\Upsilon(B)\), put
\[
  A_{\mathfrak n}
  =
  A\widehat{\otimes}_B B_{\mathfrak n}.
\]
By \SGARefLink{sga3:VIIB:0.4}{0.4}, \(A\) is a pseudocompact \(B\)-module, and hence is the product
of the \(A_{\mathfrak n}\) \textup{(\SGARefLink{sga3:VIIB:0.3.6}{0.3.6})}. The morphism \(g\) is
therefore the product of the morphisms
\[
  g_{\mathfrak n}:B_{\mathfrak n}\longrightarrow A_{\mathfrak n}
\]
obtained from \(g\) by base change. These remain epimorphisms. We are
thus reduced to the case in which \(B\) is local, with maximal ideal
\(\mathfrak n\); set \(K=B/\mathfrak n\).

By Nakayama's \SGARefLink{sga3:VIIB:0.3.3}{lemma 0.3.3}, it is enough to prove that
\[
  g\widehat{\otimes}_B K
\]
is surjective. This map is obtained from \(g\) by base change and is
therefore an epimorphism in \(\operatorname{Alp}/K\). We may thus
suppose that \(B=K\) is a field. Now \(f\) is a monomorphism if and
only if the diagonal
\[
  X\longrightarrow X\times_YX
\]
is an isomorphism; equivalently, the homomorphism
\[
  A\widehat{\otimes}_K A\longrightarrow A,
  \qquad
  x\widehat{\otimes}_Kx'\longmapsto xx',
\]
is an isomorphism. Since \(K\) is a field, this implies \(A=K\).

\medskip
\noindent\textbf{Remark.}
\phantomsection
\label{sga3:VIIB:1.3.remark}
{\itshape It follows from the proposition that every monomorphism
\(f:X\to Y\) of formal varieties is an isomorphism from \(X\) onto a
necessarily closed formal subvariety of \(Y\).}\footnote{The editors
added this remark.}

\subsubsection*{1.3.1}
\phantomsection
\label{sga3:VIIB:1.3.1}
\addcontentsline{toc}{subsubsection}{\SGARefLink{sga3:VIIA:1.3.1}{1.3.1}. Effective epimorphisms}

The preceding proposition implies, in particular, that every
monomorphism \(f:X\to Y\) in \(\operatorname{Vaf}/k\) is effective
\textup{(cf.\ \SGARefLink{sga3:iv:definition:1:3}{IV 1.3})}.\footnote{In the opposite category
\((\operatorname{Vaf}/k)^\circ=\operatorname{Alp}/k\), the
corresponding map \(g:B\to A\) is an effective epimorphism: since it is
surjective, it induces an isomorphism of profinite \(k\)-algebras
\(B/I\xrightarrow{\sim}A\), where \(I=\ker g\). Hence every morphism
\(\phi:B\to C\) in \(\operatorname{Alp}/k\) that vanishes on \(I\)
descends to a morphism \(\bar\phi:A\to C\) such that
\(\phi=\bar\phi\circ g\).}
The same is not true of epimorphisms, as one sees by slightly modifying
the counterexample of Expos\'e V, \S~2.c).\footnote{Let \(k\) be a
field, \(X=\operatorname{Spf}(k[[T]])\), and
\(Y=\operatorname{Spf}(B)\), where \(B\) is the \(k\)-subalgebra of
\(A=k[[T]]\) topologically generated by \(T^3\) and \(T^4\); thus
\(B\) consists of the series \(\sum a_nT^n\) with
\(a_n=0\) for \(n=1,2,5\). Then \(X\to Y\) is an epimorphism that is
not effective: the cokernel of
\(X\times_YX\rightrightarrows X\) is
\(\operatorname{Spf}(B')\), where
\[
  B'=\{a\in A\mid
       a\widehat{\otimes}1=1\widehat{\otimes}a\},
\]
and \(B'\) contains \(T^5\).}
This is why we now consider a useful class of effective epimorphisms.

\medskip
\noindent\textbf{Lemma.}
\phantomsection
\label{sga3:VIIB:1.3.1.flatness-criterion}
\footnote{The editors added this lemma, which explains the term
``topologically flat.''}
{\itshape Let \(f:X\to Y\) be a morphism of formal \(k\)-varieties,
let \(A\) and \(B\) be the affine algebras of \(X\) and \(Y\), and let
\(f^\natural:B\to A\) be the morphism induced by \(f\). The following
conditions are equivalent:}
\begin{enumerate}
  \item[\textup{(i)}]
  {\itshape For every \(x\in X\), the homomorphism
  \[
    f_x^\natural:
    \mathcal O_{Y,f(x)}
    \longrightarrow
    \mathcal O_{X,x}
  \]
  makes \(\mathcal O_{X,x}\) a topologically free pseudocompact
  \(\mathcal O_{Y,f(x)}\)-module.}

  \item[\textup{(ii)}]
  {\itshape For every \(y\in Y\), the local component
  \[
    A_y=\prod_{f(x)=y}\mathcal O_{X,x}
  \]
  is a topologically free pseudocompact \(B_y\)-module.}

  \item[\textup{(iii)}]
  {\itshape The map \(f^\natural:B\to A\) makes \(A\) a projective
  pseudocompact \(B\)-module.}

  \item[\textup{(iv)}]
  {\itshape The functor
  \[
    \operatorname{PC}(B)\longrightarrow\operatorname{PC}(A),
    \qquad
    M\longmapsto M\widehat{\otimes}_B A,
  \]
  is exact.}
\end{enumerate}
{\itshape When these conditions hold, \(f\) is said to be
\emph{topologically flat}.}

The implications
\[
  \textup{(i)}\Longrightarrow\textup{(ii)}
  \Longleftrightarrow\textup{(iii)}
  \Longleftrightarrow\textup{(iv)}
\]
follow from \SGARefLink{sga3:VIIB:0.2.F}{0.2.F}\,\textup{\SGARefLink{sga3:VIIB:0.2.F:item:iii}{(iii)}} and \SGARefLink{sga3:VIIB:0.3.7}{0.3.7}. Conversely, suppose that
\textup{\SGARefLink{sga3:VIIB:0.2.F:item:ii}{(ii)}} holds, and let \(x\in X\), \(y=f(x)\). Since
\(\mathcal O_{X,x}\) is a direct factor of \(A_y\), it is a projective
pseudocompact \(B_y\)-module, and therefore is topologically free by
\SGARefLink{sga3:VIIB:0.2.1}{0.2.1}, because \(B_y\) is local.

A morphism \(f:X\to Y\) of formal \(k\)-varieties is called
\emph{surjective} if it induces a surjection of the underlying sets.

\medskip
\noindent\textbf{Proposition.}
\phantomsection
\label{sga3:VIIB:1.3.1.effective-epi}
{\itshape Let \(f:X\to Y\) be a surjective, topologically flat
morphism of formal \(k\)-varieties. Then \(f\) is an effective
epimorphism \textup{(cf.\ \SGARefLink{sga3:iv:definition:1:3}{IV 1.3})}.}

Let \(A\) and \(B\) be the affine algebras of \(X\) and \(Y\), and let
\(g:B\to A\) be the morphism induced by \(f\). We must prove that \(Y\)
identifies with the cokernel of
\[
  X\times_YX\rightrightarrows X;
\]
that is, for every \(\mathfrak n\in\Upsilon(B)\),
\(B_{\mathfrak n}\) must identify with the subring of
\[
  A_{\mathfrak n}
  =
  A\widehat{\otimes}_B B_{\mathfrak n}
\]
formed by the elements \(a\) such that
\[
  a\widehat{\otimes}1=1\widehat{\otimes}a.
\]

We may suppose that \(B\) is local, with maximal ideal
\(\mathfrak n\). The hypothesis then says that \(g\) makes \(A\) a
nonzero topologically free \(B\)-module. By Nakayama's \SGARefLink{sga3:VIIB:0.3.3}{lemma 0.3.3},
\(A/\overline{\mathfrak nA}\ne0\), and hence the induced morphism
\[
  g':B/\mathfrak n\longrightarrow
     A/\overline{\mathfrak nA}
\]
is injective. By \SGARefLink{sga3:VIIB:1.3.2}{Lemma 1.3.2} below, \(B\) is a direct factor of \(A\)
as a \(B\)-module; write \(A=B\oplus A'\). It follows that \(B\)
identifies with the set of \(a\in A\) such that
\[
  a\widehat{\otimes}_B1
  =
  1\widehat{\otimes}_B a.
\]

\subsubsection*{1.3.2}
\phantomsection
\label{sga3:VIIB:1.3.2}

\noindent\textbf{Lemma 1.3.2.}
{\itshape Let \(B\) be a local pseudocompact ring, let
\(\mathfrak n\) be its maximal ideal, let \(M,N\) be projective
pseudocompact \(B\)-modules, and let \(g:M\to N\) be a morphism. If
\[
  (B/\mathfrak n)\widehat{\otimes}_B g
\]
is injective, then \(g\) is an isomorphism from \(M\) onto a direct
factor of \(N\).}

Indeed, suppose
\[
  g'=(B/\mathfrak n)\widehat{\otimes}_B g
\]
is injective. Since \(B/\mathfrak n\) is a field, \(g'\) has a
retraction \(r'\). Let \(p\) and \(q\) be the canonical projections
from \(M\) and \(N\) onto
\[
  M/\overline{\mathfrak nM}
  \qquad\text{and}\qquad
  N/\overline{\mathfrak nN}.
\]
Since \(N\) is projective, there exists \(r:N\to M\) such that
\[
  p\circ r=r'\circ q;
\]
thus \(r'\) is obtained from \(r\) by passage to the quotient. Since
\(r'\circ g'\) is an isomorphism, \SGARefLink{sga3:VIIB:0.3.4}{0.3.4} shows that the same is true of
\(r\circ g\), because \(M\) is projective. If \(s\) denotes the
inverse of \(r\circ g\), then \(s\circ r\) is a retraction of \(g\).

\subsubsection*{1.3.3}
\phantomsection
\label{sga3:VIIB:1.3.3}

\noindent\textbf{Proposition 1.3.3.}
{\itshape Let \(f:X\to Y\) and \(g:Y\to Z\) be morphisms of formal
\(k\)-varieties.}
\begin{enumerate}
  \item[\textup{(i)}]
  {\itshape If \(f\) and \(g\) are topologically flat, so is
  \(g\circ f\).}

  \item[\textup{\SGAInlineRefTarget{sga3:VIIB:1.3.3:item:ii}(ii)}]
  {\itshape If \(f\) and \(g\circ f\) are topologically flat and
  \(f\) is surjective, then \(g\) is topologically flat.}

  \item[\textup{\SGAInlineRefTarget{sga3:VIIB:1.3.3:item:iii}(iii)}]
  {\itshape If \(f\) is topologically flat, then so is
  \(f\times_YY'\) after every base change \(Y'\to Y\).}
\end{enumerate}

Assertions \textup{\SGARefLink{sga3:vii:paragraph:5-1:equation:i}{(i)}} and \textup{\SGARefLink{sga3:VIIB:1.3.3:item:iii}{(iii)}} are clear. To prove
\textup{\SGARefLink{sga3:VIIB:1.3.3:item:ii}{(ii)}}, let \(A,B,C\) be the affine algebras of \(X,Y,Z\), and
let \(f':B\to A\), \(g':C\to B\) be the homomorphisms induced by
\(f,g\). Since \(g\circ f\) is topologically flat,
\(f'\circ g'\) makes \(A\) a projective pseudocompact \(C\)-module.
Likewise, \(f'\) makes \(A\) a faithful projective pseudocompact
\(B\)-module. Thus, as \(P\) ranges over pseudocompact \(C\)-modules
and \(N\) over pseudocompact \(B\)-modules, the functors
\[
  P\longmapsto P\widehat{\otimes}_C A,
  \qquad
  N\longmapsto N\widehat{\otimes}_B A
\]
are exact; the second is also faithful. Hence the functor
\[
  P\longmapsto P\widehat{\otimes}_C B
\]
is exact, and \SGARefLink{sga3:VIIB:0.3.7}{0.3.7} implies that \(B\) is a projective
pseudocompact \(C\)-module.

\subsubsection*{1.3.4}
\phantomsection
\label{sga3:VIIB:1.3.4}

\noindent\textbf{Proposition 1.3.4.}
{\itshape Let \(S\) be a scheme, let \(Y\) be a locally Noetherian
\(S\)-scheme, and let \(f:X\to Y\) be a faithfully flat
\(S\)-morphism locally of finite type. Thus \(f\) is an effective
epimorphism; equivalently, the following sequence is exact
\textup{(cf.\ \SGARefLink{sga3:iv:proposition:6:3:1}{IV 6.3.1}\,\textup{\SGARefLink{sga3:iv:proposition:6:3:1:item:iv:01}{(iv)}} and \SGARefLink{sga3:iv:definition:1:3}{IV 1.3})}:}

\begin{center}
  \(
(\ast)\qquad
X\times_YX
\mathrel{\substack{\xrightarrow{\ \operatorname{pr}_1\ }\\[-.7ex]
                   \xrightarrow[\ \operatorname{pr}_2\ ]{}}}
X\xrightarrow{\ f\ }Y.
\)

\end{center}

{\itshape Then the morphism of formal
\(\widehat S\)-varieties
\[
  \widehat f:
  \widehat X/\widehat S
  \longrightarrow
  \widehat Y/\widehat S
\]
\textup{(cf.\ \SGARefLink{sga3:VIIB:1.2.6}{1.2.6})} is surjective and topologically flat, and the
following sequence, obtained from the preceding one, is exact:}

\begin{center}
  \(
\widehat{(\ast)}\qquad
\widehat{X\times_YX}/\widehat S
\mathrel{\substack{\xrightarrow{\ \widehat{\operatorname{pr}}_1\ }\\[-.7ex]
                   \xrightarrow[\ \widehat{\operatorname{pr}}_2\ ]{}}}
\widehat X/\widehat S
\xrightarrow{\ \widehat f\ }
\widehat Y/\widehat S.
\)

\end{center}

Let \(y\) be a point of \(Y\), with image \(s\) in \(S\), such that
\(\kappa(y)\) is a finite extension of the residue field
\(\kappa(s)\). Since \(f\) is surjective and locally of finite type,
\(f^{-1}(y)\) is nonempty and locally of finite type over
\(\kappa(y)\). Its closed points are precisely the points of
\(\widehat X/\widehat S\) mapping to \(y\). This proves that
\(\widehat f\) is surjective.

Let \(x\) be a closed point of
\(f^{-1}(y)\).\footnote{The editors expanded the argument that follows
and then used the material added in \SGARefLink{sga3:VIIB:1.2.6}{1.2.6}.} Since \(Y\) is locally
Noetherian and \(f\) is locally of finite type, the local rings
\(\mathcal O_{Y,y}\) and \(\mathcal O_{X,x}\) are Noetherian. The
powers of their maximal ideals therefore have finite colength, and the
local rings of \(\widehat Y/\widehat S\) at \(y\) and of
\(\widehat X/\widehat S\) at \(x\) are the completions
\[
  \widehat{\mathcal O}_{Y,y}
  \qquad\text{and}\qquad
  \widehat{\mathcal O}_{X,x}
\]
for their adic topologies. Since \(f\) is flat,
\(\widehat{\mathcal O}_{X,x}\) is flat over
\(\widehat{\mathcal O}_{Y,y}\) by SGA~1, IV 5.8. By \SGARefLink{sga3:VIIB:0.3.8}{0.3.8} it is
therefore a topologically free
\(\widehat{\mathcal O}_{Y,y}\)-module. Thus \(\widehat f\) is
topologically flat.

\SGARefLink{sga3:VIIB:1.3.1}{Proposition 1.3.1} now shows that \(\widehat f\) is an effective
epimorphism. In other words, the following sequence is exact, where
\(\widehat X\) is written for \(\widehat X/\widehat S\):

\begin{center}
  \(
\widehat X\times_{\widehat Y}\widehat X
\mathrel{\substack{\xrightarrow{\ \widehat{\operatorname{pr}}_1\ }\\[-.7ex]
                   \xrightarrow[\ \widehat{\operatorname{pr}}_2\ ]{}}}
\widehat X\xrightarrow{\ \widehat f\ }\widehat Y.
\)

\end{center}

Moreover, \SGARefLink{sga3:VIIB:1.2.6}{1.2.6} gives a natural isomorphism, compatible in particular
with the projections onto \(\widehat X\):

\begin{center}
  \(
\widehat{X\times_YX}
\simeq
\widehat X\times_{\widehat Y}\widehat X.
\)

\end{center}

Consequently, the completed sequence displayed in the proposition is
exact.

